\section{Introduction}

Consensus clustering was first studied by Ailon, Charikar and Newman \cite{ailon2008aggregating}. Given a set of input clusterings 
\[
\m{I} = \{ \m{C}_1, \m{C}_2, \ldots, \m{C}_m \}
\]
defined on a common set of vertices $V$, the goal is to compute a clustering $\m{C}^*$ that minimizes the aggregated pairwise distance to all input clusterings. 
Formally, we seek
\[
    \m{C}^* = \argmin_{\m{C}} \left( \sum_{i = 1}^m \dist(\m{C}_i, \m{C}) \right),
\]
where $\dist(\m{C}_i, \m{C}_j)$ denotes the pairwise disagreement between clusterings $\m{C}_i$ and $\m{C}_j$, measured as the number of vertex pairs that are clustered together in one clustering but separated in the other:
\[
    \dist(\m{C}_i, \m{C}_j) = 
    \big| \{ (u,v) \mid 
    (u \sim_{\m{C}_i} v \land u \nsim_{\m{C}_j} v)
    \lor 
    (u \nsim_{\m{C}_i} v \land u \sim_{\m{C}_j} v)
    \} \big|.
\]
Here, $u \sim_{\m{C}} v$ indicates that vertices $u$ and $v$ belong to the same cluster in $\m{C}$.

\subsection{Our Contributions}

\begin{restatable}{theorem}{lazymedoid}[High probability upper bound]
\label{thm:sampled-lazy-medoid}
    Let \(\varepsilon,\eta\in(0,1]\) and \(\delta\in(0,1)\). There is a randomized online algorithm for consensus clustering that, with probability at least \(1-\delta\), satisfies for every time step \(t\),
\[
    \obj_t(S_t)
    \le
    (2+\varepsilon)(1+\eta)\obj_t(\OPT_t).
\]
Moreover, on the same event, every vertex has recourse at most
\[
    O\left(
        \frac{1}{\varepsilon\eta}
        \log\frac{n}{\delta}
        \log(mn)
    \right).
\]
\end{restatable}

\begin{restatable}{theorem}{zerorecoupper}[In expectation zero recourse upper bound]
    \label{thm:random-zero-recourse}
    There is a randomized online algorithm for consensus clustering such that, for every time step \(t\),
    \[
        \mathbb{E}\bigl[\obj_t(S_t)\bigr] \le 2\left(1-\frac1m\right)\obj_t(\OPT_t).
    \]
    Moreover, the algorithm has zero recourse.
\end{restatable}

\begin{restatable}[Deterministic zero-recourse lower bound]{theorem}{detzerorecourselb}
\label{thm:det-zero-recourse-lb}
For every deterministic zero-recourse online algorithm and every integer \(q\ge 1\), there exists an online consensus clustering instance with \(q+2\) vertices and \(4\) input clusterings such that, at the final time \(T\),
\[
    \obj_T(S_T)
    \ge
    \left(2-\frac{1}{q+1}\right)\obj_T(\OPT_T).
\]
Consequently, every deterministic zero-recourse online algorithm has competitive ratio at least
\[
    2-o(1).
\]
\end{restatable}

\begin{restatable}[Randomized zero-recourse lower bound]{theorem}{randzerorecourselb}
\label{thm:rand-zero-recourse-lb}
For every randomized zero-recourse online algorithm and every integer \(q\ge 1\), there exists an online consensus clustering instance with \(q+2\) vertices and \(4\) input clusterings such that, at the final time \(T\),
\[
    \mathbb E[\obj_T(S_T)]
    \ge
    \left(
        \frac32-\frac{1}{2(q+1)}
    \right)\obj_T(\OPT_T),
\]
where the expectation is over the algorithm's internal randomness. Consequently, every randomized zero-recourse online algorithm has expected competitive ratio at least
\[
    \frac32-o(1).
\]
\end{restatable}

\kushagra{Look at the preliminaries for the definition of notations.}
